
\chapter{Estado del arte}
\label{cap:estado-arte}

En este capítulo se revisan los trabajos y tecnologías relacionados con el proyecto. Las referencias se guardan en \texttt{bibliografia/bibliografia.bib} y se citan con \verb|\cite{clave}|. Solo aparecen en la bibliografía las entradas citadas.

\section{Ejemplos de citas}

Un libro clásico sobre algoritmos es el de Knuth~\cite{knuth1997}. Dijkstra argumentó en contra de la sentencia \texttt{goto}~\cite{dijkstra1968}, y la arquitectura Transformer se presentó en un congreso~\cite{vaswani2017}. Se pueden citar varias referencias a la vez~\cite{knuth1997,dijkstra1968,vaswani2017}, así como documentación en línea~\cite{python-docs}, software~\cite{latex-project} o trabajos académicos~\cite{apellido2025}.

Para las citas textuales se usa \verb|\enquote{...}|, que pone las comillas adecuadas al idioma: \enquote{la simplicidad es un requisito previo de la fiabilidad}.
